\subsection*{Matched multi-model intervention (E5)}

The controlled results above establish that the seven applicability invariants
are executable and causally separable under oracle and predicted structured
states (E1--E4). They do not by themselves establish that supplying an explicit
situation state improves a frontier language model on natural prompts. That is
the E5 claim in the evidence ladder, and it requires model responses rather than
symbolic controls. We therefore evaluate several proprietary and open-weight
models on a stratified, evidence-complete rendering of SituationCatch-Bench
(30 items per category), holding the evidence block, decoding temperature and
output-token budget fixed and varying only the prompt condition. The primary
quantity is the situation-on-versus-off contrast within each model: the
situation-state condition minus the strongest non-situational prompt for the
same model, so any difference is attributable to explicit state rather than to
model identity, evidence or budget. Raw responses, usage, latencies and failed
calls are archived per model; the reported accuracies read directly from those
records and exclude failed calls, which are reported as $n$.

% E5 multi-model table -- PENDING.
% No real model rows were found in the summary CSV yet (only mock/empty data).
% Run:  python code/run_e5.py        (needs OPENROUTER_API_KEY with credits)
% then: python code/make_e5_table.py  to regenerate this file with real numbers.
\begin{table}[t]
\caption{E5 multi-model situation-intervention results \emph{(pending execution;
no frontier-model responses are included in this build)}. Regenerated
automatically from \texttt{paper/results/multimodel\_summary.csv} once
\texttt{code/run\_e5.py} has produced real responses.}
\label{tab:e5-multimodel}
\centering\small
\begin{tabular}{l}
\toprule
Awaiting E5 execution --- see \texttt{code/experiment\_manifest.e5.json}. \\
\bottomrule
\end{tabular}
\end{table}


We interpret a positive within-model $\Delta$ as support for the hypothesis that
explicit situation state reduces applicability errors under matched conditions,
not as evidence of deployment readiness: the rendered prompts still derive from
the seven generated scenario families, so natural-text and distributional
generalization (E5 splits and E6) remain open. The frozen study manifest,
provider gateway, model snapshots and cost accounting are recorded in
\texttt{code/experiment\_manifest.e5.json}, and the run is reproduced by a single
command (\texttt{python code/run\_e5.py}).
